\documentclass[11pt,a4paper]{article}
\newcommand{\intestazione}{Dimero di spin --- 6. VQE sotto rumore}
% =====================================================================
%  Preambolo comune ai documenti sul dimero (serie "che cosa abbiamo fatto")
%  Incluso con % =====================================================================
%  Preambolo comune ai documenti sul dimero (serie "che cosa abbiamo fatto")
%  Incluso con % =====================================================================
%  Preambolo comune ai documenti sul dimero (serie "che cosa abbiamo fatto")
%  Incluso con \input{preambolo_dimero} subito dopo \documentclass.
% =====================================================================
\usepackage[T1]{fontenc}
\usepackage[utf8]{inputenc}
\usepackage[main=italian,provide=*]{babel}
\usepackage{amsmath,amssymb,mathtools}
\usepackage{graphicx}
\usepackage{booktabs}
\usepackage{colortbl}
\usepackage{xcolor}
\usepackage{geometry}
\usepackage{placeins}
\usepackage{caption}
\usepackage{enumitem}
\usepackage{fancyhdr}
\usepackage{needspace}
\usepackage{tikz}
\usetikzlibrary{arrows.meta,positioning,shapes.geometric,calc,fit,backgrounds}
\usepackage{hyperref}
\hypersetup{colorlinks=true,linkcolor=blu,citecolor=blu,urlcolor=blu}

\geometry{margin=2.5cm}

% --- colori coerenti con quelli delle figure (palette Okabe-Ito)
\definecolor{blu}{HTML}{0072B2}
\definecolor{arancio}{HTML}{E69F00}
\definecolor{verdes}{HTML}{009E73}
\definecolor{rossos}{HTML}{D55E00}
\definecolor{violas}{HTML}{CC79A7}
\definecolor{grigetto}{HTML}{E8E8E8}

\captionsetup{font=small,labelfont=bf,width=0.92\textwidth}

\pagestyle{fancy}
\fancyhf{}
\fancyhead[L]{\small\itshape\intestazione}
\fancyhead[R]{\small\thepage}
\renewcommand{\headrulewidth}{0.3pt}

% --- notazione
\newcommand{\ket}[1]{\lvert #1\rangle}
\newcommand{\bra}[1]{\langle #1 \rvert}
\newcommand{\media}[1]{\langle #1 \rangle}
\newcommand{\vs}[1]{\mathbf{s}_{#1}}

% --- riquadro "in una riga": il messaggio da portare a casa
\newcommand{\messaggio}[1]{%
  \begin{center}
  \begin{minipage}{0.93\textwidth}
  \setlength{\fboxsep}{9pt}%
  \colorbox{grigetto}{\begin{minipage}{\dimexpr\textwidth-2\fboxsep}
  \small\textbf{In una riga.}\enspace #1
  \end{minipage}}
  \end{minipage}
  \end{center}
}

% --- riquadro "come lo sappiamo": la verifica numerica dietro un'affermazione
\newcommand{\verifica}[1]{%
  \begin{center}
  \begin{minipage}{0.93\textwidth}
  \small\textcolor{verdes}{\rule{2.2em}{1.1pt}}\enspace
  \textcolor{verdes}{\textbf{Come lo sappiamo.}}\enspace #1
  \end{minipage}
  \end{center}
  \medskip
}
 subito dopo \documentclass.
% =====================================================================
\usepackage[T1]{fontenc}
\usepackage[utf8]{inputenc}
\usepackage[main=italian,provide=*]{babel}
\usepackage{amsmath,amssymb,mathtools}
\usepackage{graphicx}
\usepackage{booktabs}
\usepackage{colortbl}
\usepackage{xcolor}
\usepackage{geometry}
\usepackage{placeins}
\usepackage{caption}
\usepackage{enumitem}
\usepackage{fancyhdr}
\usepackage{needspace}
\usepackage{tikz}
\usetikzlibrary{arrows.meta,positioning,shapes.geometric,calc,fit,backgrounds}
\usepackage{hyperref}
\hypersetup{colorlinks=true,linkcolor=blu,citecolor=blu,urlcolor=blu}

\geometry{margin=2.5cm}

% --- colori coerenti con quelli delle figure (palette Okabe-Ito)
\definecolor{blu}{HTML}{0072B2}
\definecolor{arancio}{HTML}{E69F00}
\definecolor{verdes}{HTML}{009E73}
\definecolor{rossos}{HTML}{D55E00}
\definecolor{violas}{HTML}{CC79A7}
\definecolor{grigetto}{HTML}{E8E8E8}

\captionsetup{font=small,labelfont=bf,width=0.92\textwidth}

\pagestyle{fancy}
\fancyhf{}
\fancyhead[L]{\small\itshape\intestazione}
\fancyhead[R]{\small\thepage}
\renewcommand{\headrulewidth}{0.3pt}

% --- notazione
\newcommand{\ket}[1]{\lvert #1\rangle}
\newcommand{\bra}[1]{\langle #1 \rvert}
\newcommand{\media}[1]{\langle #1 \rangle}
\newcommand{\vs}[1]{\mathbf{s}_{#1}}

% --- riquadro "in una riga": il messaggio da portare a casa
\newcommand{\messaggio}[1]{%
  \begin{center}
  \begin{minipage}{0.93\textwidth}
  \setlength{\fboxsep}{9pt}%
  \colorbox{grigetto}{\begin{minipage}{\dimexpr\textwidth-2\fboxsep}
  \small\textbf{In una riga.}\enspace #1
  \end{minipage}}
  \end{minipage}
  \end{center}
}

% --- riquadro "come lo sappiamo": la verifica numerica dietro un'affermazione
\newcommand{\verifica}[1]{%
  \begin{center}
  \begin{minipage}{0.93\textwidth}
  \small\textcolor{verdes}{\rule{2.2em}{1.1pt}}\enspace
  \textcolor{verdes}{\textbf{Come lo sappiamo.}}\enspace #1
  \end{minipage}
  \end{center}
  \medskip
}
 subito dopo \documentclass.
% =====================================================================
\usepackage[T1]{fontenc}
\usepackage[utf8]{inputenc}
\usepackage[main=italian,provide=*]{babel}
\usepackage{amsmath,amssymb,mathtools}
\usepackage{graphicx}
\usepackage{booktabs}
\usepackage{colortbl}
\usepackage{xcolor}
\usepackage{geometry}
\usepackage{placeins}
\usepackage{caption}
\usepackage{enumitem}
\usepackage{fancyhdr}
\usepackage{needspace}
\usepackage{tikz}
\usetikzlibrary{arrows.meta,positioning,shapes.geometric,calc,fit,backgrounds}
\usepackage{hyperref}
\hypersetup{colorlinks=true,linkcolor=blu,citecolor=blu,urlcolor=blu}

\geometry{margin=2.5cm}

% --- colori coerenti con quelli delle figure (palette Okabe-Ito)
\definecolor{blu}{HTML}{0072B2}
\definecolor{arancio}{HTML}{E69F00}
\definecolor{verdes}{HTML}{009E73}
\definecolor{rossos}{HTML}{D55E00}
\definecolor{violas}{HTML}{CC79A7}
\definecolor{grigetto}{HTML}{E8E8E8}

\captionsetup{font=small,labelfont=bf,width=0.92\textwidth}

\pagestyle{fancy}
\fancyhf{}
\fancyhead[L]{\small\itshape\intestazione}
\fancyhead[R]{\small\thepage}
\renewcommand{\headrulewidth}{0.3pt}

% --- notazione
\newcommand{\ket}[1]{\lvert #1\rangle}
\newcommand{\bra}[1]{\langle #1 \rvert}
\newcommand{\media}[1]{\langle #1 \rangle}
\newcommand{\vs}[1]{\mathbf{s}_{#1}}

% --- riquadro "in una riga": il messaggio da portare a casa
\newcommand{\messaggio}[1]{%
  \begin{center}
  \begin{minipage}{0.93\textwidth}
  \setlength{\fboxsep}{9pt}%
  \colorbox{grigetto}{\begin{minipage}{\dimexpr\textwidth-2\fboxsep}
  \small\textbf{In una riga.}\enspace #1
  \end{minipage}}
  \end{minipage}
  \end{center}
}

% --- riquadro "come lo sappiamo": la verifica numerica dietro un'affermazione
\newcommand{\verifica}[1]{%
  \begin{center}
  \begin{minipage}{0.93\textwidth}
  \small\textcolor{verdes}{\rule{2.2em}{1.1pt}}\enspace
  \textcolor{verdes}{\textbf{Come lo sappiamo.}}\enspace #1
  \end{minipage}
  \end{center}
  \medskip
}

\newcommand{\Tr}{\operatorname{Tr}}

\title{\bfseries Un risultato ideale, eseguito su un simulatore imperfetto\\[2pt]
\large Documento 6 --- VQE con termine DM sotto rumore}
\author{Samuele Schiavi \\ \small Tesi triennale in Fisica, Università di Parma
--- Relatore: Prof.\ Alessandro Chiesa}
\date{}

\begin{document}
\maketitle
\thispagestyle{fancy}

\section{La prima domanda della Parte 2}

Con il modello di rumore pronto (Documento 5), la prima cosa naturale da
chiedere è la più semplice della catena: il VQE con termine DM del Documento
2 --- gli stessi parametri $\theta^*$, lo stesso ansatz PMA-2q$\cdot$3, lo
stesso punto di lavoro ``test 2'' ($b/J=0.35$, $D/J=0.80$) --- trova ancora
uno stato buono se il circuito gira su un simulatore rumoroso invece che a
vettore di stato esatto?

Una scelta di modellazione va dichiarata subito, perché guida come leggere
tutto il resto di questo documento: si \emph{riusano} i parametri già
ottimizzati nel caso ideale, senza rioptimizzare vedendo il rumore durante
l'ottimizzazione. Questo documento misura quindi la degradazione di un
risultato ideale, non la convergenza di un VQE che ottimizzi già sotto
rumore --- quella è una domanda diversa, affrontata a parte (Documento 10).

\section{Energia e fedeltà, versione rumorosa}

Con il circuito rumoroso lo stato preparato non è più puro: è un operatore
densità $\rho$. L'energia si generalizza in modo diretto,
$E_\text{rumoroso}=\Tr[\rho H]$. La fedeltà rispetto allo stato di riferimento
esatto $\ket{\psi_\text{ideale}}$ (fedeltà di Uhlmann, semplificata quando uno
dei due stati è puro) diventa
\begin{equation}
F(\rho,\ket{\psi_\text{ideale}}) = \bra{\psi_\text{ideale}}\rho\ket{\psi_\text{ideale}},
\end{equation}
verificata contro la routine standard di Qiskit su uno stato misto casuale ---
non il caso specifico di questo problema, apposta per evitare una verifica
circolare.

\section{Risultato}

\begin{center}
\renewcommand{\arraystretch}{1.3}
\begin{tabular}{@{}lccc@{}}
\toprule
\textbf{Condizione} & $E$ & $F$ & CNOT\\
\midrule
Rumore nullo (verifica di coerenza) & $-3.57321145$ & $1.0000000000$ & $2$\\
Rumore di riferimento (\texttt{ibm\_torino}) & $-3.50164605$ & $0.9850060074$ & $2$\\
\bottomrule
\end{tabular}
\end{center}

Il caso a rumore nullo riproduce l'energia esatta della Parte 1 a precisione
numerica: un primo controllo di coerenza interna, superato prima di guardare
al caso rumoroso. Al rumore di riferimento, lo scostamento di energia è
$+0.071565$ e la perdita di fedeltà $1-F=0.014994$, poco sotto l'1.5\%.

\verifica{Il caso a rumore nullo ripete a precisione macchina l'energia
$E_0=-3.57321145$ già trovata nel Documento 2 --- il primo controllo che va
sempre fatto prima di fidarsi di un numero rumoroso: se il ramo ``rumore
spento'' non torna, non ha senso guardare il ramo acceso.}

\messaggio{Un ansatz a soli 2 CNOT lascia poco spazio al rumore di
accumularsi: $1.5\%$ di fedeltà persa è la scala di riferimento più piccola
di tutta la pipeline --- gli stadi successivi, con circuiti più lunghi,
partiranno da qui e la degraderanno ulteriormente.}

\section{Come degrada al variare del rumore di gate}

\begin{figure}[htbp]
\centering
\includegraphics[width=0.92\textwidth]{figure/fig_passo2_scan_eps2q}
\caption{Energia e fedeltà rumorose in funzione di $\varepsilon_{2q}$, a
$\varepsilon_{1q}$ e $p_\text{readout}$ fissati al valore di riferimento. La
linea tratteggiata segna il punto di calibrazione \texttt{ibm\_torino}.
Degradazione monotona e liscia, nessuna soglia o comportamento anomalo
nell'intervallo esplorato --- non ci si aspetta un massimo o un minimo qui,
a differenza di quanto succederà con la dinamica temporale.}
\end{figure}

\FloatBarrier

\section{Che cosa non dice ancora questo documento}

Un solo punto (o uno scan di un solo parametro) fissa la scala della
degradazione per uno stato \emph{statico}: non dice come si comporti la
fedeltà se il circuito viene allungato --- qui non c'è alcuna dinamica --- né
se il rumore possa essere assorbito da un'ottimizzazione diversa dei
parametri. La prima domanda è quella naturale del passo successivo
(Documento 7); la seconda è affrontata a parte (Documento 10).

\messaggio{I parametri $\theta^*$ e il loro costo minimo (2 CNOT, solo per la
preparazione) sono il punto di partenza del documento successivo, che
aggiunge sopra questa stessa preparazione $N$ passi di evoluzione temporale.
La scala di degradazione qui misurata ($1-F\sim1.5\%$) sarà lì il contributo
costante, indipendente da $N$, dentro una fedeltà totale che invece dipende
da $N$ in modo tutt'altro che monotono.}

\end{document}
